\documentclass[12pt]{article}
\usepackage[a4paper, portrait, margin=0.8in]{geometry}
\usepackage{graphicx}

\author{Aurélien Casteilla}
\date{26-04-2021}
\title{Mansion64}

\begin{document}
\pagenumbering{gobble}
\maketitle
\begin{figure}[h!]
    \includegraphics[width=\linewidth]{mascotte.png}
\end{figure}
\newpage

\paragraph{}
        Mansion64 is a 64 bits little endian ISA. This ISA tries to apply
        the KISS principle and also tries to respect the Popek and Goldberg
        virtualization requirement. Also, this CPU will include atomic memory
        operation for a multi-core version. So Mansion64 is a RISC ISA.
        Mansion64 has a 64 bits-wide data bus, a 64 bits-wide address bus,
        two sets of 31 64 bits-wide registers for a total of 47 registers, plus
        a zero register.  There is also a status register and a saved status
        register. One set of register is used by the user's program and the
        other set is reserved for fast interrupt handling or for the supervisor
        usage. The instruction word is 32 bit wide.
        \setcounter{secnumdepth}{-1}
\section{Summary}
\begin{itemize}
    \item RISC ISA
    \item applies the KISS principle
    \item respects the Popek and Goldberg virtualization requirement
    \item have atomic memory operations
    \item 64 bits or less address bus (more than you will ever need (16 EiB))
    \item 64 bits data bus
    \item 64 bits registers
    \item 32 bits instruction word
    \item 15 registers shared between the user and the supervisor
    \item 16 registers are user specific and 16 registers are supervisor specific 
    \item 32 registers visible at the same time (including a zero register)
    \item one address space for everything (Von Neumann architecture with I/O mapped in memory)
    \item "user" and "supervisor" mode
    \item one instruction per clock cycle if the architecture is pipelined    
	   (excluding memory instructions especially if the memory is slow)
    \item SIMD instructions on 2 words, 4 half-words or 8 bytes
    \item free and open source
\end{itemize}
\section{Notes :}
\begin{itemize}
    \item First, this ISA will be implemented on an emulator that runs on Linux
        (and maybe other *nix). This ISA is intended to be implemented on
        hardware the first time on a FPGA development board (Arduino MKR Vidor
        4000, with an Intel Cyclone 10) The final goal is to make a
        mini-computer implementation of this ISA.

    \item I am aware that Intel Quartus is not free software.

    \item This is just an ISA for the CPU part of a computer. So it doesn't
        have and doesn't specify a MMU, a FPU, a GPU, a cache, a PIC, or an I/O
        controller. However, some of these elements will be on the FPGA SoC.

    \item I named it "Mansion64" because a mini-computer is really big and you
        will need a castle or a mansion for it (don't worry, I'm kidding). 64
        because it's a 64 bits architecture. Also "Mansion" because this is a
        freak and uncanny architecture.

\end{itemize}
\section{Glossary of the acronyms :}

\paragraph{CPU : }Central Processing Unit
\paragraph{FPGA : }Field Programmable Gate Array
\paragraph{FPU : }Floating Point Unit
\paragraph{FOSS : }Free and Open Sourced Software
\paragraph{GPU : }Graphical Processing Unit
\paragraph{I/O : }Input / Output
\paragraph{ISA : }Instruction Set Architecture
\paragraph{KISS : }Keep It Simple Stupid
\paragraph{MMU : }Memory Management Unit
\paragraph{PIC : }Programmable Interrupt Controller
\paragraph{RISC : }Reduced Instruction Set Computing
\paragraph{SIMD : }Single Instruction Multiple Data 
\paragraph{SoC : }System On Chip

\section{}
Cyclone 10 is a trademark owned by Intel

\end{document}
